\documentclass[11pt]{article}
\usepackage[utf8]{inputenc}
\usepackage[T1]{fontenc}
\usepackage[spanish,es-nodecimaldot]{babel}
\usepackage[a4paper,margin=2.5cm]{geometry}
\usepackage{amsmath,amssymb,booktabs,graphicx,float,hyperref}
\newcommand{\figroot}{../../outputs/toi_ati_case_anatomy/figures_pdf}
\newcommand{\safeincludegraphics}[2][]{%
  \IfFileExists{#2}{\includegraphics[#1]{#2}}{\fbox{\begin{minipage}{0.86\textwidth}\centering Figura no encontrada: \texttt{#2}\end{minipage}}}%
}
\title{Anatomia de casos TOI/ATI: regiones Mapper y planetas ancla}
\author{Proyecto Mapper/TDA de exoplanetas}
\date{}
\begin{document}
\maketitle

\section{Introduccion}
Este reporte abre el ranking global TOI/ATI para explicar por que ciertas regiones Mapper y planetas ancla fueron priorizados. El objetivo no es afirmar descubrimientos de planetas, sino auditar la anatomia de los indices.

\section{Que problema resuelve la anatomia TOI/ATI}
La salida principal es una anatomia de los casos top: que factor empuja cada indice, que factor lo limita y que tan estable parece el deficit relativo al cambiar de radio.

\section{Datos usados}
El modulo consume tablas generadas por \texttt{topological\_incompleteness\_index} y, cuando existen, tablas contextuales de sombra observacional y estudios locales.

\section{Descomposicion de TOI}
\[
\mathrm{TOI}(v)=\mathrm{Shadow}(v)(1-I_{R^3}(v))C_{\mathrm{phys}}(v)S_{\mathrm{net}}(v).
\]
TOI debe leerse como una prioridad regional que combina frontera observacional, baja imputacion, continuidad fisica con vecinos y soporte de red.

\section{Descomposicion de ATI}
\[
\mathrm{ATI}(p^*)=\mathrm{TOI}(v)\Delta_{\mathrm{rel,best}}(p^*)(1-I_{R^3}(p^*))A(p^*).
\]
ATI debe leerse como una prioridad local de inspeccion. El termino $\Delta_{\mathrm{rel,best}}$ debe revisarse junto con los radios individuales.

\section{Top regiones}
\begin{figure}[H]\centering
\safeincludegraphics[width=.9\textwidth]{\figroot/top_regions_toi_case_anatomy.pdf}
\caption{Top regiones por TOI.}
\end{figure}

\begin{figure}[H]\centering
\safeincludegraphics[width=.9\textwidth]{\figroot/toi_factor_decomposition.pdf}
\caption{Descomposicion visual de factores TOI.}
\end{figure}

\section{Top planetas ancla}
\begin{figure}[H]\centering
\safeincludegraphics[width=.9\textwidth]{\figroot/top_anchors_ati_case_anatomy.pdf}
\caption{Top planetas ancla por ATI.}
\end{figure}

\begin{figure}[H]\centering
\safeincludegraphics[width=.9\textwidth]{\figroot/ati_factor_decomposition.pdf}
\caption{Descomposicion visual de factores ATI.}
\end{figure}

\section{Deficit por radio}
\begin{figure}[H]\centering
\safeincludegraphics[width=.8\textwidth]{\figroot/deficit_by_radius_summary.pdf}
\caption{Resumen del deficit relativo por radio.}
\end{figure}

\section{Discusion}
Una region con TOI alto puede ganar por sombra fuerte, baja imputacion, continuidad fisica o soporte de red. Un ancla con ATI alto puede ganar por pertenecer a una region TOI alta, por mostrar deficit local, por baja imputacion o por ser representativa del nodo.

\section{Limitaciones}
Este modulo no introduce una funcion de completitud instrumental, no estima cantidades absolutas de objetos reales ausentes y no convierte por si solo un caso priorizado en una conclusion observacional cerrada.

\section{Conclusion}
Este modulo no calcula otro indice nuevo; abre los rankings TOI/ATI para explicar por que ciertas regiones y planetas ancla fueron priorizados. La salida principal es una anatomia de los casos top, no una afirmacion de objetos ausentes confirmados.
\end{document}
